\section{Metodologia}
\label{sec:metodologia}

\subsection{Visão Geral do Processo de Avaliação}
\label{subsec:visao-geral-metodologia}

O método de avaliação conduzido neste trabalho fundamenta-se no processo de referência estabelecido pela norma ISO/IEC 25040:2011 \cite{iso25040_2011}. Para o escopo da presente entrega (Release E1), a execução compreende as duas primeiras fases do processo normativo: o estabelecimento dos requisitos da avaliação (Fase 1) e a especificação da avaliação (Fase 2), detalhadas passo a passo nas Subseções~\ref{sec:fase1} e \ref{sec:fase2-gqm}. A abordagem metodológica caracteriza-se como uma avaliação empírica de produto baseada em observação externa (caixa-preta), cujo objetivo é comparar o Portal de Serviços e o aplicativo móvel oficial do Detran-DF quanto às características de adequação funcional, eficiência de desempenho e cobertura de contexto \cite{iso25010_2023,iso25010_2011}.

Para assegurar rigor analítico e rastreabilidade estrita entre os requisitos das partes interessadas e a coleta de evidências, a especificação na Fase 2 orienta-se pela abordagem GQM \cite{basili1994gqm}, desdobrando objetivos conceituais em questões analíticas, hipóteses falseáveis, métricas quantitativas, níveis de pontuação e critérios de julgamento. Paralelamente, a governança interna da equipe apoia-se no modelo PSM/CID \cite{mcgarry2002psm}, com rastreabilidade de esforço, progresso e qualidade de revisão documentada por meio de \textit{issues} e \textit{pull requests} no repositório oficial do projeto.